\documentclass[10pt,letterpaper]{article}
\usepackage{cloudsuite}
\renewcommand{\DocShort}{CH. 02}
\begin{document}
\coverpage{Chapter 2: Case Study Background}{Cornell study notes and applied review | ORIENTATION}{Concepts \textbullet{} decisions \textbullet{} verification artifacts}
\sourceStatement
\tableofcontents
\newpage
\section{Learning contract}
\begin{takeaway}
Analyze organizations through infrastructure context, business drivers and migration strategy.
\end{takeaway}
\textbf{Chapter scope:} ATN, DTGOV and Innovartus Technologies Inc..\\
\sourceRef{2}{11}
\subsection{Mapped source sections}
\begin{itemize}
\item 2.1 Case Study \#1: ATN
\item 2.2 Case Study \#2: DTGOV
\item 2.3 Case Study \#3: Innovartus Technologies Inc.
\end{itemize}
\section{Cornell cue and notes matrix}
\small\setlength{\LTleft}{0pt}\setlength{\LTright}{0pt}
\begin{longtable}{>{\RaggedRight\arraybackslash}p{0.29\textwidth}>{\RaggedRight\arraybackslash}p{0.63\textwidth}}
\toprule\textbf{Cue / essential question}&\textbf{Explanatory notes and reasoning}\\\midrule\endfirsthead\toprule\textbf{Cue}&\textbf{Notes}\\\midrule\endhead
\textbf{Why introduce three case studies?} & ATN, DTGOV and Innovartus supply distinct organizational contexts used to illustrate later concepts. Their constraints should not be merged into a fictional single environment.\\[3pt]\textcolor{petrol}{\emph{Practice:} Make a separate context sheet for each named case.}\\[8pt]\midrule
\textbf{What belongs in infrastructure context?} & Capture workload inventory, legacy dependencies, network and data boundaries, operational tooling and technical debt. This is the baseline against which cloud changes can be compared.\\[3pt]\textcolor{petrol}{\emph{Practice:} List five existing dependencies before proposing migration.}\\[8pt]\midrule
\textbf{What belongs in a business driver?} & Distinguish cost reduction, agility, scalability and new service capabilities. A business goal requires a measure, not merely an aspiration.\\[3pt]\textcolor{petrol}{\emph{Practice:} Convert a goal into a measurable service or cost outcome.}\\[8pt]\midrule
\textbf{How should a roadmap be justified?} & Stage the implementation against dependencies, verification and organizational readiness. Distinguish preparatory migration from actual service availability.\\[3pt]\textcolor{petrol}{\emph{Practice:} Prepare a milestone, evidence and risk table.}\\[8pt]\midrule
\textbf{How can a case study support architecture learning?} & Treat scenarios as decision exercises: state the constraint, choose an applicable mechanism or architecture, and identify what would falsify the choice.\\[3pt]\textcolor{petrol}{\emph{Practice:} Write one architectural hypothesis for each organization.}\\[8pt]\midrule
\end{longtable}\normalsize
\section{Synthesis and architecture connections}
Translate the chapter vocabulary into one explicit service requirement, one boundary, one owner and one measurable test. Later mechanism and architecture chapters provide possible realizations; the requirements remain the acceptance baseline.
\begin{workbox}[title={Applied deliverable}]
Analyze organizations through infrastructure context, business drivers and migration strategy.\\[4pt]Create a one-page architecture or operations artifact with: (1) assumptions, (2) a decision and alternatives, (3) control ownership, (4) observable acceptance criteria, and (5) a failed-case test.
\end{workbox}
\section{Active recall}
\begin{enumerate}
\item Why introduce three case studies?
\item What belongs in infrastructure context?
\item What belongs in a business driver?
\item How should a roadmap be justified?
\end{enumerate}
\subsection{Answer-check rubric}
\begin{itemize}
\item A correct response distinguishes the named concepts rather than relying on vendor names.
\item The explanation identifies an operational consequence, boundary or tradeoff.
\item At least one measurement, test or documentary artifact could demonstrate the claim.
\end{itemize}
\section{Blank Cornell capture sheet}
\textbf{Topic / source subsection:}\enspace\rule{0.65\linewidth}{0.4pt}\\[6pt]
\begin{longtable}{p{0.28\textwidth}p{0.64\textwidth}}
\toprule Cues / questions & Notes / evidence \\ \midrule
\rule{0pt}{42pt} & \\ \midrule
\rule{0pt}{42pt} & \\ \midrule
\rule{0pt}{42pt} & \\ \midrule
\rule{0pt}{42pt} & \\ \midrule
\rule{0pt}{42pt} & \\ \midrule
\end{longtable}
\textbf{Summary in three sentences:}\\[16pt]\hrulefill\\[16pt]\hrulefill\\[16pt]\hrulefill

\refs
\end{document}
